\documentclass[11pt]{article}
\usepackage[margin=1in]{geometry}
\usepackage{amsmath,amssymb,amsthm,mathtools}
\usepackage{lmodern}
\usepackage{booktabs}
\usepackage{enumitem}
\usepackage{xcolor}
\usepackage[hidelinks]{hyperref}
\hypersetup{
  pdftitle={A Canonical Additive Edge Frame for Zero-Hole Reduced-Residue BDH Remainders},
  pdfauthor={Liang Wang},
  pdfsubject={Complete-graph tight frame, exact diagonal deletion, and literal-edge no-sparsification}
}

\newtheorem{theorem}{Theorem}[section]
\newtheorem{proposition}[theorem]{Proposition}
\newtheorem{lemma}[theorem]{Lemma}
\newtheorem{corollary}[theorem]{Corollary}
\newtheorem{remark}[theorem]{Remark}
\newtheorem{definition}[theorem]{Definition}

\newcommand{\Fq}{\mathbb F_q}
\newcommand{\C}{\mathbb C}
\newcommand{\e}{\mathrm e}
\newcommand{\one}{\mathbf 1}
\newcommand{\cQ}{\mathcal Q}
\newcommand{\cE}{\mathcal E}
\newcommand{\cK}{\mathcal K}
\newcommand{\supp}{\operatorname{supp}}
\newcommand{\rank}{\operatorname{rank}}
\newcommand{\diag}{\operatorname{diag}}

\title{\textbf{A Canonical Additive Edge Frame for Zero-Hole\\
Reduced-Residue BDH Remainders}}
\author{Liang Wang\\
Huazhong University of Science and Technology\\
Wuhan 430074, P.R. China}
\date{August 2026}

\begin{document}
\emergencystretch=2em
\maketitle

\begin{abstract}
We identify the exact additive-frequency geometry of the standard-zero-hole
reduced-residue Barban--Davenport--Halberstam remainder.  For a prime modulus
$q$, the zero-hole variance is $q^{-1}y^*P_{q-1}y$, where $y$ contains the
$q-1$ nonzero additive Fourier coefficients and $P_{q-1}$ projects off their
constant direction.  The complete-graph Laplacian turns this projection into
a tight frame of $(q-1)(q-2)/2$ literal two-frequency differences.  The same
edge cells distribute the mandatory $(q-2)/(q-1)$ coefficient diagonal
exactly, leaving a pure coefficient-off-diagonal remainder in every cell.
Four-packet complex polarization is then edgewise, and contraction of the
frame recovers the original physical residue kernel with no normalization
change.  We also prove a scoped uniqueness theorem: every scalar-weighted
representation by literal vectors $e_k-e_l$ must use every edge with weight
$1/(q-1)$, so no strict edge subset represents the projection.  A
residue-zero spike refutes separate absolute estimates of equal and unequal
frequencies: the two pieces are nonzero, opposite, and sum to zero.  An
independent exact certificate checks 431 finite identities for
$q=2,3,5,7,11$.  The result is structural: it creates an additive
pre-emitter, not a Kloosterman attachment or a prime-shell power saving.
\end{abstract}

\section{Introduction and claim firewall}

The prime-dynamics Gate-B reduction studied in the preceding stages produces
a signed, prime-modulus, kernel-localized reduced-residue variance.  Complex
polarization converts its two-sequence scalar into four one-sequence
remainders, but the coefficient diagonal must be subtracted with the exact
factor $q-2$ before the packet signs are reassembled
\cite{Wang2026V59}.  Translating a physical block changes the distinguished
deleted residue; the moving-hole correction has already been isolated and
paid at the critical $1/96$ clock \cite{Wang2026TPC207}.  The remaining
arithmetic object is therefore the standard-zero-hole remainder itself.

A first additive-Fourier attempt would split that remainder into equal- and
unequal-frequency terms and estimate them separately.  This order is invalid.
For a row supported only at residue zero, the true zero-hole variance is zero,
whereas the two proposed pieces have equal nonzero magnitude and opposite
sign.  Taking absolute values before their cancellation manufactures a large
term from an exact zero.

The stable object is the complete-graph Laplacian on the nonzero additive
frequencies.  It retains the constant-direction cancellation internally and
simultaneously distributes the mandatory coefficient diagonal.  The paper's
contributions are the following.

\begin{enumerate}[leftmargin=2.2em]
  \item We prove that the zero-hole variance is $q^{-1}y^*P_{q-1}y$ and give
  its exact complete-graph tight-frame expansion.  The frame has rank $q-2$,
  $(q-1)(q-2)/2$ edge vectors, and redundancy $(q-1)/2$ for $q>2$.
  \item We prove the edge-mass identity
  $\sum_e|\Delta_e(n)|^2=q(q-2)\mathbf 1_{q\nmid n}$.  It places the
  $(q-2)/(q-1)$ coefficient diagonal inside the same edge cells, so every
  emitted cell is coefficient-off-diagonal before any estimate.
  \item We lift the four-packet polarization edge by edge, recover the frozen
  physical residue kernel exactly, and expose an oriented $(d,k)$ difference
  fiber with the unit-annihilating factor $1-\e_q(-dn)$.
  \item We prove that a scalar-weighted decomposition by literal
  two-frequency differences is unique.  Every edge is forced, which rules out
  edge-subset sparsification and identifies the collective source theorem
  that is still missing.
\end{enumerate}

All four items are finite-dimensional identities proved for every prime
$q$ and every finitely supported complex sequence.  They do not estimate the
edge cells.  In particular, the paper does not prove a prime-only
Barban--Davenport--Halberstam asymptotic, attach the cells to a Kloosterman
theorem, pay the full strict $1/400$ Gate-B margin, establish an arithmetic
$L^2$ statement, earn fixed-atom credit, or prove a twin-prime theorem.

\section{The frozen zero-hole row}

Let $q$ be prime and put
\[
 \e_q(t)=\exp(2\pi i t/q),\qquad \e(t)=\exp(2\pi i t).
\]
Let $a=(a_n)$ be a finitely supported complex sequence, let $H>0$, and let
$v\in\mathbb R$.  Write
\begin{equation}\label{eq:residue-row}
 x_n(v)=a_n\e(vn/H),\qquad
 A_r(a;v)=\sum_{n\equiv r\; (\mathrm{mod}\ q)}x_n(v).
\end{equation}
The mean over the $q-1$ retained residues and the standard-zero-hole variance
are
\begin{equation}\label{eq:zero-hole-variance}
 \overline A^{\times}=\frac1{q-1}\sum_{r\ne0}A_r,
 \qquad
 V_0(a;v)=\sum_{r\ne0}|A_r-\overline A^{\times}|^2.
\end{equation}
The mandatory coefficient diagonal and its remainder are
\begin{equation}\label{eq:diagonal-remainder}
 D_0(a)=\frac{q-2}{q-1}\sum_{q\nmid n}|a_n|^2,
 \qquad R_0(a;v)=V_0(a;v)-D_0(a).
\end{equation}
The physical Gate-B row carries the outer factor $qR_0$, not merely $R_0$.

We use the additive transform
\begin{equation}\label{eq:dft}
 \widehat A(k)=\sum_{r\bmod q}A_r\e_q(-kr)
 =\sum_n a_n\e(vn/H)\e_q(-kn).
\end{equation}
This convention is fixed throughout.  No multiplicative character, inverse
residue, or asymptotic estimate enters the frame construction.

\section{Zero-hole projection in additive frequency}

Set $d=q-1$, enumerate the nonzero frequencies, and define
\begin{equation}\label{eq:projection}
 y=(\widehat A(k))_{k\in\Fq^\times}\in\C^d,
 \qquad P_d=I_d-\frac1d\one\one^*.
\end{equation}

\begin{theorem}[zero-hole frequency projection]\label{thm:projection}
For every prime $q$, every finitely supported complex sequence $a$, and every
real $v$,
\begin{equation}\label{eq:projection-identity}
 \boxed{V_0(a;v)=\frac1q\,y^*P_dy.}
\end{equation}
Moreover, $P_d$ is an orthogonal projection of rank $q-2$.  For $q=2$ both
sides of \eqref{eq:projection-identity} vanish.
\end{theorem}

\begin{proof}
For a general row $z=(z_r)_{r\bmod q}$, let
$\mu=q^{-1}\sum_rz_r$ and let $V_{\rm all}=\sum_r|z_r-\mu|^2$.
Deleting coordinate zero and recentering the retained coordinates gives the
elementary leave-one-out identity
\begin{equation}\label{eq:leave-one-out}
 \sum_{r\ne0}\left|z_r-\frac1{q-1}\sum_{s\ne0}z_s\right|^2
 =V_{\rm all}-\frac q{q-1}|z_0-\mu|^2.
\end{equation}
Indeed, the retained mean equals
$\mu-(z_0-\mu)/(q-1)$; expanding around $\mu$ proves
\eqref{eq:leave-one-out}.

Apply this identity to $z_r=A_r$.  Parseval and Fourier inversion yield
\begin{equation}\label{eq:parseval-inversion}
 V_{\rm all}=\frac1q\sum_{k\ne0}|\widehat A(k)|^2,
 \qquad
 A_0-\mu=\frac1q\sum_{k\ne0}\widehat A(k).
\end{equation}
Substitution gives
\begin{equation}\label{eq:projection-expanded}
 V_0=\frac1q\sum_{k\ne0}|\widehat A(k)|^2
 -\frac1{q(q-1)}\left|\sum_{k\ne0}\widehat A(k)\right|^2,
\end{equation}
which is \eqref{eq:projection-identity}.  The matrix $P_d$ projects onto
$\one^\perp$, so its rank is $d-1=q-2$.  When $q=2$, $d=1$ and $P_d=0$.
\end{proof}

The null direction in Theorem~\ref{thm:projection} is essential.  A vector
whose nonzero Fourier coordinates are all equal contributes no zero-hole
variance, even if its Euclidean norm is large.

\section{The complete-graph tight frame}

Let $K_d$ be the complete graph on the nonzero-frequency vertices.  For an
unordered edge $e=\{k,l\}$ define
\begin{equation}\label{eq:edge-vector}
 g_e=e_k-e_l,
 \qquad
 \Delta_{k,l}(n)=\e_q(-kn)-\e_q(-ln).
\end{equation}
The graph Laplacian is
\begin{equation}\label{eq:laplacian}
 \sum_{e\in E(K_d)}g_eg_e^*=dI_d-\one\one^*=dP_d.
\end{equation}

\begin{theorem}[canonical additive edge frame]\label{thm:edge-frame}
For every object in Theorem~\ref{thm:projection},
\begin{equation}\label{eq:edge-frame}
 \boxed{
 V_0(a;v)=\frac1{q(q-1)}
 \sum_{\{k,l\}\in E(K_{q-1})}
 \left|\sum_n a_n\e(vn/H)\Delta_{k,l}(n)\right|^2.}
\end{equation}
The edge family contains $(q-1)(q-2)/2$ vectors and spans a space of
dimension $q-2$.  For $q>2$ its frame redundancy is $(q-1)/2$.
\end{theorem}

\begin{proof}
The edge transform is
\begin{equation}\label{eq:edge-transform}
 T_e[a](v)=\langle g_e,y\rangle
 =\widehat A(k)-\widehat A(l)
 =\sum_na_n\e(vn/H)\Delta_{k,l}(n).
\end{equation}
Contracting \eqref{eq:laplacian} with $y$ gives
$\sum_e|T_e[a](v)|^2=d\,y^*P_dy$.  Theorem~\ref{thm:projection} and
$d=q-1$ prove \eqref{eq:edge-frame}.  The count and redundancy follow from
the complete graph and the rank of $P_d$.  For $q=2$ the edge set is empty.
\end{proof}

The redundancy in Theorem~\ref{thm:edge-frame} is exact bookkeeping.  It is
not a power saving and does not permit a triangle inequality over the edge
set.  Section~\ref{sec:uniqueness} will show that every literal edge is forced
in a scalar-weighted representation of this projection.

\section{Exact edgewise diagonal deletion}

Each edge difference annihilates the excluded residue:
\begin{equation}\label{eq:zero-residue-annihilation}
 q\mid n\quad\Longrightarrow\quad \Delta_{k,l}(n)=0.
\end{equation}
The total edge mass on a unit residue is also exact.

\begin{lemma}[edge mass]\label{lem:edge-mass}
For every integer $n$,
\begin{equation}\label{eq:edge-mass}
 \boxed{
 \sum_{e\in E(K_{q-1})}|\Delta_e(n)|^2
 =q(q-2)\mathbf 1_{q\nmid n}.}
\end{equation}
\end{lemma}

\begin{proof}
If $q\mid n$, \eqref{eq:zero-residue-annihilation} proves the claim.  Suppose
$q\nmid n$ and put
$u(n)=(\e_q(-kn))_{k\ne0}$.  Then
$\|u(n)\|_2^2=q-1$ and
$\langle\one,u(n)\rangle=\sum_{k\ne0}\e_q(-kn)=-1$.  Therefore
\begin{align*}
 \sum_e|\langle g_e,u(n)\rangle|^2
 &=u(n)^*\bigl((q-1)I-\one\one^*\bigr)u(n)\\
 &=(q-1)^2-1=q(q-2).
\end{align*}
\end{proof}

Define the edge diagonal and edge remainder by
\begin{equation}\label{eq:edge-diagonal}
 D_e[a]=\sum_n|a_n|^2|\Delta_e(n)|^2,
 \qquad
 \cE_e^\circ[a](v)=|T_e[a](v)|^2-D_e[a].
\end{equation}
Expanding the square removes the coefficient diagonal inside each cell:
\begin{equation}\label{eq:pure-off-diagonal}
 \boxed{
 \cE_e^\circ[a](v)=
 \sum_{t\ne u}a_t\overline{a_u}\e(v(t-u)/H)
 \Delta_e(t)\overline{\Delta_e(u)}.}
\end{equation}

\begin{corollary}[zero-hole off-diagonal edge compiler]
\label{cor:edge-compiler}
The diagonal in \eqref{eq:diagonal-remainder} distributes over the edge frame
without an error:
\begin{equation}\label{eq:diagonal-distribution}
 \frac1{q(q-1)}\sum_eD_e[a]
 =\frac{q-2}{q-1}\sum_{q\nmid n}|a_n|^2=D_0(a).
\end{equation}
Consequently,
\begin{equation}\label{eq:edge-remainder}
 \boxed{
 R_0(a;v)=\frac1{q(q-1)}\sum_e\cE_e^\circ[a](v),
 \qquad
 qR_0(a;v)=\frac1{q-1}\sum_e\cE_e^\circ[a](v).}
\end{equation}
\end{corollary}

\begin{proof}
Sum \eqref{eq:edge-diagonal} over edges and apply
Lemma~\ref{lem:edge-mass}.  Subtract the result from
Theorem~\ref{thm:edge-frame}.  The second identity retains the mandatory
physical outer factor $q$.
\end{proof}

Corollary~\ref{cor:edge-compiler} is the main local improvement over a bare
Fourier decomposition.  The frame geometry and coefficient-diagonal
cancellation now occupy the same cell; neither has to be recovered after an
absolute estimate.

\section{Polarization and the physical kernel}

For two finitely supported sequences $\beta,w$, define the bilinear edge cell
\begin{equation}\label{eq:bilinear-edge-cell}
 \cE_e^\circ(\beta,w;v)
 =T_e[\beta](v)\overline{T_e[w](v)}
 -\sum_n\beta(n)\overline{w(n)}|\Delta_e(n)|^2.
\end{equation}
It has the pure off-diagonal expansion
\begin{equation}\label{eq:bilinear-off-diagonal}
 \cE_e^\circ(\beta,w;v)
 =\sum_{t\ne u}\beta(t)\overline{w(u)}\e(v(t-u)/H)
 \Delta_e(t)\overline{\Delta_e(u)}.
\end{equation}

\begin{proposition}[edgewise four-packet polarization]
\label{prop:polarization}
Let $a^{(j)}=\beta+i^jw$ for $j=0,1,2,3$.  Then, for every edge,
\begin{equation}\label{eq:edge-polarization}
 \boxed{
 \frac14\sum_{j=0}^3i^j\cE_e^\circ[a^{(j)}](v)
 =\cE_e^\circ(\beta,w;v).}
\end{equation}
\end{proposition}

\begin{proof}
For complex $x,y$,
$x\overline y=\frac14\sum_{j=0}^3i^j|x+i^jy|^2$.
Apply this identity first to $T_e[\beta],T_e[w]$ and then pointwise to
$\beta(n),w(n)$.  Subtracting the two polarized identities gives
\eqref{eq:edge-polarization}.
\end{proof}

Let $\psi_+$, the prime shell $\cQ$, and the literal sequences $\beta,w$ be
the frozen V59 objects.  Corollary~\ref{cor:edge-compiler}, summed before any
outer absolute value, gives the exact physical normal form
\begin{equation}\label{eq:physical-scalar}
 \boxed{
 \mathfrak C_x=
 \int_{\mathbb R}\psi_+(v)
 \sum_{q\in\cQ}\frac1{q-1}
 \sum_{e\in E(K_{q-1})}\cE_e^\circ(\beta,w;v)\,dv.}
\end{equation}
The ordered block decomposition also commutes with
\eqref{eq:physical-scalar}: replace $\beta,w$ by each ordered pair
$\beta_b,w_c$ and sum over $(b,c)$ before applying an absolute value.  If the
$v$-integral is evaluated first, its kernel $K_H$ produces the additive
pre-emitter
\begin{equation}\label{eq:pre-emitter}
 \sum_{t\ne u}\beta_b(t)\overline{w_c(u)}K_H(u-t)
 \Delta_e(t)\overline{\Delta_e(u)}.
\end{equation}
Equation~\eqref{eq:pre-emitter} is not yet a Kloosterman cell.

The following contraction verifies every residue coefficient in
\eqref{eq:physical-scalar}.  For $r,s\bmod q$, put
\begin{equation}\label{eq:physical-kernel-definition}
 \cK_q(r,s)=\sum_e
 (\e_q(-kr)-\e_q(-lr))(\e_q(ks)-\e_q(ls)).
\end{equation}

\begin{proposition}[physical-kernel crosswalk]\label{prop:physical-kernel}
The edge contraction is
\begin{equation}\label{eq:physical-kernel}
 \boxed{
 \cK_q(r,s)=
 \begin{cases}
 0,&rs\equiv0\pmod q,\\
 q(q-2),&r=s\ne0,\\
 -q,&r,s\ne0,\ r\ne s.
 \end{cases}}
\end{equation}
Hence, on unit residues,
\begin{equation}\label{eq:literal-crosswalk}
 \frac{\cK_q(r,s)}{q-1}
 =q\left(\mathbf1_{r=s}-\frac1{q-1}\right),
\end{equation}
which is the original V59 coefficient.
\end{proposition}

\begin{proof}
If either residue is zero, every edge difference on that coordinate vanishes.
For nonzero $r,s$, use \eqref{eq:laplacian} with
$u(r)=(\e_q(-kr))_{k\ne0}$.  The diagonal case is
$(q-1)^2-1=q(q-2)$ as in Lemma~\ref{lem:edge-mass}.  If $r\ne s$, additive
orthogonality gives $\langle u(s),u(r)\rangle=-1$, while both coordinate sums
equal $-1$; contraction with $(q-1)I-\one\one^*$ therefore gives
$(q-1)(-1)-1=-q$.  Dividing by $q-1$ proves
\eqref{eq:literal-crosswalk}.
\end{proof}

\section{Oriented difference fibers and uniqueness}
\label{sec:uniqueness}

Write $l=k+d$ in $\Fq$.  Counting each unordered edge in both orientations
gives
\begin{equation}\label{eq:oriented-sum}
 \sum_{e\in E(K_{q-1})}\cE_e^\circ
 =\frac12\sum_{d\in\Fq^\times}
 \sum_{\substack{k\in\Fq^\times\\k\ne-d}}
 \cE_{k,k+d}^\circ.
\end{equation}
The edge factor separates as
\begin{equation}\label{eq:difference-fiber}
 \boxed{
 \Delta_{k,k+d}(n)=\e_q(-kn)\bigl(1-\e_q(-dn)\bigr).}
\end{equation}
Thus \eqref{eq:physical-scalar} is equivalently
\begin{equation}\label{eq:oriented-physical-scalar}
 \mathfrak C_x=
 \int\psi_+(v)\sum_{q\in\cQ}\frac1{2(q-1)}
 \sum_{d\ne0}\sum_{k\ne0,-d}
 \cE_{k,k+d}^\circ(\beta,w;v)\,dv.
\end{equation}
The base twist $k$ and the unit-annihilating difference factor in
\eqref{eq:difference-fiber} are the preferred interface for a future
Poisson/Kloosterman compiler.  The factor $1/2$ and the exclusion
$k=-d$ are mandatory.

The complete edge family cannot be reduced within its literal class.

\begin{theorem}[literal-edge no-sparsification]\label{thm:no-sparsification}
Let $d\ge1$.  Suppose scalar weights $w_{k,l}\in\C$ satisfy
\begin{equation}\label{eq:weighted-edge-decomposition}
 P_d=\sum_{1\le k<l\le d}
 w_{k,l}(e_k-e_l)(e_k-e_l)^*.
\end{equation}
Then
\begin{equation}\label{eq:forced-edge-weight}
 \boxed{w_{k,l}=\frac1d=\frac1{q-1}}
\end{equation}
for every edge.  In particular, no strict subset of literal
two-frequency differences represents $P_d$ with scalar edge weights.
\end{theorem}

\begin{proof}
For $k\ne l$, the $(k,l)$ entry of $P_d$ is $-1/d$.  On the right side of
\eqref{eq:weighted-edge-decomposition}, only the outer product belonging to
$\{k,l\}$ has a nonzero $(k,l)$ entry, and that entry is $-w_{k,l}$.  Hence
$w_{k,l}=1/d$ for every pair.  With these weights the diagonal entries agree
automatically by the complete-graph identity.
\end{proof}

The scope in Theorem~\ref{thm:no-sparsification} matters.  It excludes neither
a dense orthonormal basis of $\one^\perp$, a signed higher-rank cell system,
nor an arithmetic theorem that estimates the full redundant frame jointly.
It says that deleting literal edges and repairing their weights cannot be the
next compiler.

\section{Sharp falsifiers and mutation boundaries}

\subsection{Equal/off-equal estimation is unstable}

Take $A_0=L$ and $A_r=0$ for every $r\ne0$.  The retained row is zero, so
$V_0=0$.  On the other hand, $\widehat A(k)=L$ for every $k\ne0$, and the two
terms in \eqref{eq:projection-expanded} are
\begin{equation}\label{eq:spike-equal}
 \frac1q\sum_{k\ne0}|\widehat A(k)|^2
 =\frac{q-1}{q}|L|^2,
\end{equation}
\begin{equation}\label{eq:spike-off-equal}
 -\frac1{q(q-1)}\left|\sum_{k\ne0}\widehat A(k)\right|^2
 =-\frac{q-1}{q}|L|^2.
\end{equation}
Separate absolute estimates lose their exact cancellation.  Every edge
difference in Theorem~\ref{thm:edge-frame} instead vanishes before an estimate
is taken.  This fixture refutes the proposed estimate order, not the valid
algebraic split in \eqref{eq:projection-expanded}.

\subsection{One coefficient cancels cellwise}

If $a$ has one nonzero coefficient at a unit residue, then
$|T_e[a]|^2=D_e[a]$ for every edge.  Thus
$\cE_e^\circ[a]=0$ edge by edge.  A global diagonal subtraction that leaves
nonzero individual cells would fail this local test.

\subsection{Mandatory factors}

Several nearby formulas describe different objects.
\begin{itemize}[leftmargin=2em]
  \item Summing oriented edges without the factor $1/2$ doubles the scalar.
  \item Including frequency zero replaces $K_{q-1}$ by $K_q$ and produces an
  all-residue projection rather than the zero-hole row.
  \item Replacing $q-1$ by $q$ changes every off-diagonal projection entry and
  violates Theorem~\ref{thm:no-sparsification}.
  \item Omitting $D_e[a]$ returns a positive variance, not the signed
  coefficient-off-diagonal remainder.
  \item Dropping one literal edge changes exactly one pair of off-diagonal
  matrix entries and cannot be repaired by reweighting other edges.
\end{itemize}

These mutations are included in the exact checker.  Their rejection protects
the physical signs and normalization; it does not supply arithmetic evidence.

\section{Source boundary and related analytic engines}

The additive edge frame is an upstream representation theorem.  Three nearby
analytic sources clarify what remains to be proved.

\paragraph{General-sequence BDH architecture.}
Harper proves simple Barban--Davenport--Halberstam type asymptotics for broad
classes of general sequences under explicit Progressions,
Non-concentration, and further structural hypotheses
\cite{Harper2024}.  For prime moduli, grouping reduced residues naturally
leads to a zero-hole variance.  The source does not state
\eqref{eq:physical-scalar}, verify its hypotheses uniformly for the four
literal prime-dynamics packets, select the prime subset after the signed
diagonal subtraction, or perform the ordered block reassembly.  The present
paper neither invokes nor disproves Harper's theorem; it gives an exact local
normal form for testing a future attachment.

\paragraph{Post-emitter Kloosterman bounds.}
Blomer and Pascadi obtain strong bounds for bilinear forms with Kloosterman
sums, including the critical fixed-modulus saving used in the route ledger
\cite{BlomerPascadi2026}.  Their theorem accepts an already emitted form with
source-valid coefficient arrays.  The additive expression
\eqref{eq:pre-emitter} has not been transformed into those arrays.  In
particular, the complete $(d,k)$ frame, four packet signs, prime shell, block
norms, and one final reassembly cannot be supplied by renaming
\eqref{eq:pre-emitter} a Kloosterman cell.

\paragraph{Exceptional-spectrum and sparse-Fourier machinery.}
Pascadi's large-sieve framework for exceptional Maass forms supplies powerful
sparse-Fourier and incomplete-Kloosterman tools after the relevant transform
and norm hypotheses have been established \cite{Pascadi2026LargeSieve}.
Here those hypotheses are precisely part of the open compiler.  The source is
a plausible downstream engine, not evidence that the edge pre-emitter already
meets its input contract.

\paragraph{Bounded novelty claim.}
The complete-graph Laplacian identity is standard finite-dimensional linear
algebra.  A bounded arXiv metadata search on 17 August 2026 found no direct
match for the combinations ``reduced residue and graph Laplacian,'' ``BDH and
additive Fourier,'' or ``leave-one-out variance and Fourier.''  Such a query
does not prove literature-wide novelty.  The claim made here is narrower: the
exact attachment to the frozen zero-hole remainder, edgewise distribution of
its mandatory diagonal, physical-kernel crosswalk, and scoped uniqueness of
its literal two-frequency representation.

\section{Exact computational certificate}

The accompanying producer works with Gaussian rationals and prime
cyclotomic coefficient vectors.  For $q=2,3,5,7,11$, it checks the
complete-graph Laplacian entrywise, reduces every physical-kernel entry
exactly, compares direct and emitted variances on deterministic complex rows,
verifies complex polarization, and exercises the spike and edge-weight
falsifiers.  The independent checker imports neither the producer nor its
mathematical module; it reconstructs the projection and closed physical
kernel from separate formulas.

The canonical certificate contains 431 exact mathematical rows:
\begin{center}
\begin{tabular}{lr}
\toprule
Check family & Exact rows\\
\midrule
Laplacian entries, edge counts, and ranks & 167\\
Physical residue kernel & 208\\
Row variance and diagonal normalization & 20\\
Complex polarization & 2\\
Falsifiers and literal-edge uniqueness & 20\\
Explicit mutation checks & 14\\
\midrule
Total & 431\\
\bottomrule
\end{tabular}
\end{center}
Both programs use explicit failures rather than optimization-sensitive
assertions, reject duplicate JSON keys and nonfinite values at their trust
boundaries, and reproduce byte-identical output in normal and optimized
Python modes.  These computations detect transcription and normalization
mutations.  The proofs in Sections 3--7, not the finite table, establish the
general identities.

\section{Route consequence and open theorem}

Table~\ref{tab:status} separates the completed structural layer from the
unpaid arithmetic layer.

\begin{table}[ht]
\centering
\caption{Claim status for the zero-hole additive edge route.}
\label{tab:status}
\begin{tabular}{p{0.53\textwidth}p{0.32\textwidth}}
\toprule
Statement & Status\\
\midrule
Zero-hole projection and complete-graph frame & \texttt{PROVED}\\
Edgewise $(q-2)$ diagonal deletion & \texttt{PROVED}\\
Physical-kernel and four-packet crosswalk & \texttt{PROVED}\\
No strict literal-edge subset representation & \texttt{PROVED}, scoped\\
Separate absolute equal/off-equal estimates & \texttt{REFUTED}\\
Whole-frame Poisson/Kloosterman emission & \texttt{OPEN}\\
Prime-shell fixed saving greater than $1/400$ & \texttt{OPEN / UNPAID}\\
Arithmetic advance, fixed atom, or $L^2$ theorem & \texttt{NONE}\\
\bottomrule
\end{tabular}
\end{table}

The smallest next theorem is now explicit.

\begin{quote}
\textbf{Open collective compiler.}
Apply Möbius and Poisson or Voronoi transformations to the complete oriented
$(d,k)$ frame in \eqref{eq:oriented-physical-scalar} before any edge or fiber
triangle inequality.  Produce source-valid Kloosterman cells and retain a
fixed saving after the ordered blocks, four packet signs, and prime moduli are
reassembled once.
\end{quote}

The most concrete heuristic is to look for a dual variable shared across the
entire tight frame.  A shared variable could preserve the complete-graph
orthogonality through the transform.  If each edge or each $d$-fiber acquires
an independent absolute loss instead, Theorem~\ref{thm:no-sparsification}
predicts a prohibitive multiplicity and the proposed compiler should be
recorded as obstructed rather than promoted conditionally.

\section{Conclusion}

The standard-zero-hole remainder has a canonical additive representation:
the complete-graph Laplacian on nonzero frequencies, with the coefficient
diagonal deleted inside each edge cell.  This representation preserves the
outer $q$ normalization, the four-packet signs, and the original physical
residue kernel.  It also explains why the earlier equal/off-equal estimate
order fails and why literal edge-subset sparsification cannot repair it.

The resulting object is a precise pre-emitter.  The arithmetic burden has not
shrunk into a local bound: it has become a joint whole-frame transformation
and reassembly theorem.  Future progress requires that theorem, or a rigorous
obstruction showing that the complete frame cannot share enough dual
structure under the available Poisson/Kloosterman engines.

\bibliographystyle{plain}
\bibliography{references}

\end{document}
